% Direct-authority diplomatic French for printed p.123.
\SourcePage{123}

\LayerTarget{tohoku:section:1.1:p4}
\noindent bijectif (la réciproque étant en général inexacte).

\LayerTarget{tohoku:section:1.1:p5}
Le composé de deux monomorphismes (épimorphismes) est un monomorphisme
(épimorphisme) donc le composé de deux bijections est une bijection ; de même
le composé de deux isomorphismes est un isomorphisme. Si le composé $vu$ de
deux morphismes $u,v$ est un monomorphisme (resp. un épimorphisme) alors $u$
(resp. $v$) l'est aussi. Bien que le développement de tels sorites soit
évidemment nécessaire, nous nous dispenserons le plus souvent par la suite de
les dérouler explicitement, et nous nous contenterons d'indiquer avec soin les
définitions.

\LayerTarget{tohoku:section:1.1:p6}
Considérons deux monomorphismes $u:B\longrightarrow A$ et
$u':B'\longrightarrow A$, on dit que \textit{$u'$ majore ou contient $u$} et
on écrit $u\leq u'$, si on peut factoriser $u$ en $u'v$, où $v$ est un
morphisme de $B$ dans $B'$ (qui est alors déterminé de façon unique). C'est là
une relation de \textit{préordre} dans la classe des monomorphismes à valeurs
dans $A$. On dira que deux tels monomorphismes $u,u'$ sont
\textit{équivalents} si chacun majore l'autre, alors les morphismes
correspondants $B\longrightarrow B'$ et $B'\longrightarrow B$ sont inverses
l'un de l'autre. Choisissons (par exemple au moyen du symbole à tout faire
$\tau$ de Hilbert) un monomorphisme dans toute classe de monomorphismes
équivalents : les monomorphismes choisis seront appelés les
\textit{sous-trucs de $A$}. Ainsi un sous-truc de $A$ est, non un simple objet
de $C$, mais un objet $B$ muni d'un monomorphisme $u:B\longrightarrow A$,
appelé \textit{l'injection canonique} de $B$ dans $A$. (Néanmoins, par abus de
langage, on désignera souvent un sous-truc de $A$. par le symbole $B$ de
l'objet $\in C$ correspondant). La relation de majoration définit une relation
\textit{d'ordre} (et non seulement de préordre) sur la classe des sous-trucs de
$A$. De ce qu'on a vu plus haut résulte que les sous-trucs de $A$ contenus dans
un sous-truc $B$ s'identifient aux sous-trucs de $B$, cette correspondance
respectant les relations d'ordre naturelles. (Ceci ne signifie pas toutesfois
qu'un sous-truc de $B$ soit \textit{égal} à un sous-truc de $A$, ce qui
exigerait en effet que l'on ait $A=B$).

\LayerTarget{tohoku:section:1.1:p7}
Dualement, la considération d'un préordre sur la classe des épimorphismes de
$A$ permet de définir la classe ordonnée des \textit{trucs quotient de $A$}.

\LayerTarget{tohoku:section:1.1:p8}
Soient $A\in C$, et soit $(u_i)_{i\in I}$ une famille non vide de morphismes
$u_i:A\longrightarrow A_i$. Alors pour tout $B\in C$, les applications
$v\longrightarrow u_iv$ de $\mathrm{Hom}(B,A)$ dans $\mathrm{Hom}(B,A_i)$
définissent une application naturelle
\[
  \mathrm{Hom}(B,A)\longrightarrow
  \prod_{i\in I}\mathrm{Hom}(B,A_i).
\]
On dit que les $u_i$ définissent une \textit{représentation de $A$ comme
produit direct des $A_i$}, si quel que soit $B$, l'application précédente est
bijective. S'il en est ainsi, et si $A'$ est un autre objet de $C$ représenté
comme produit des $A'_i$ par des morphismes
$u'_i:A'\longrightarrow A'_i$, (l'ensemble d'indices étant le même), alors
pour toute famille $(v_i)$ de morphismes $v_i:A_i\longrightarrow A'_i$ il
existe un morphisme et un seul $v$ de $A$ dans $A'$ tel que
$u'_iv=v_iu_i$ pour tout $i$. On en conclut que si les $v_i$ sont des
équivalences, il en est de même de $v$ ; en particulier, si les $v_i$ sont les
applications identiques $I_{A_i}$, on voit que deux objects $A,A'$ représentés
comme produits de la famille des $A_i$ sont canoniquement isomorphes. Il est
donc

\TohokuContinuation{124}
